\documentclass[10pt,a4paper]{amsart}
\usepackage[margin=19mm]{geometry}
\usepackage{amsmath,amssymb,mathtools}
\usepackage[scr=rsfs,cal=euler]{mathalfa}
\usepackage{xfrac}
\usepackage[unicode,hidelinks,bookmarks=false]{hyperref}
\newcommand{\ie}{i.e.\ }
\newcommand{\cf}{cf.\ }
\newcommand{\resp}{respectively\ }
\newcommand{\Subsection}{\subsection}
\providecommand{\nobreakdash}{\nobreak\hbox{-}}
\setlength{\emergencystretch}{2em}
\multlinegap0pt
\usepackage{amssymb}
\usepackage{enumerate}
\newtheorem{theorem}{Theorem}
\renewcommand{\leq}{\leqslant}
\title{Elementary amenable groups of cohomological dimension 3}
\author{Jonathan A. Hillman}
\date{Source: arXiv:2102.02947v3 (CC BY 4.0)}
\begin{document}
\maketitle

\noindent\emph{Source and licence.} Elementary amenable groups of cohomological dimension 3. Version arXiv:2102.02947v3 (CC BY 4.0), distributed by arXiv under the Creative Commons Attribution 4.0 International licence, http://creativecommons.org/licenses/by/4.0/, as recorded in the arXiv licence metadata for this exact version and shown on the abstract page https://arxiv.org/abs/2102.02947v3. Full licence notice: \texttt{licenses/arxiv-2102.02947-CC-BY-4.0.txt}.\par\smallskip

\noindent\textit{Editorial scope.} Counted unit: the article's theorem h=3 (source0.tex lines 195--203). The article's complete original proof of it is delivered with it (source0.tex lines 205--248, one \texttt{proof} environment of 44 lines). No mathematics is written, repaired or re-derived: the statement and the proof are the article's own lines, in the article's own notation, and no retained line is substituted or re-flowed.  Retained uncounted as the source-local context the proof needs: lines 112--118.  Nothing was omitted from any retained span, so the package carries no omission record.  No retained line contains a citation, so the wrapper carries no bibliography and no bibliography entry.  No retained line contains a figure environment, an image-inclusion command or drawing code, so no image is delivered and the package declares no asset dependency.  Editorial formatting only, in addition: an AMS \texttt{amsart} preamble that restates the article's own declarations and macros used by the retained lines, namely the environments \texttt{theorem} on separate counters, together with the standard packages those lines need.  The retained environments print in this document as Theorem 1 in order of appearance, whereas the article prints the same items with the numbering of its own full document.  The complete native source of the article as distributed in the arXiv e-print for this exact version is bundled unchanged as \texttt{sources/107779/source0.tex}.
\begin{theorem}
\label{h=2}
Let $G$ be a torsion-free elementary amenable group of Hirsch length $2$.
Then $\sqrt{G}$ is abelian,
and either $\sqrt{G}$ has rank $1$ and $G\cong\sqrt{G}\rtimes\mathbb{Z}$
or $\sqrt{G}$ has rank $2$ and $[G:\sqrt{G}]\leq2$.
\end{theorem}

\begin{theorem}
\label{h=3}
Let $G$ be a torsion-free elementary amenable group of Hirsch length $3$.
If $h(\sqrt{G})=1$ then $\sqrt{G}$ is abelian and $G/\sqrt{G}\cong\mathbb{Z}^2$.
If $h(\sqrt{G})=2$ then $\sqrt{G}$ is abelian and $G/\sqrt{G}\cong\mathbb{Z}$, 
$D_\infty$ or $\mathbb{Z}\oplus\mathbb{Z}/2\mathbb{Z}$.
If $h(\sqrt{G})=3$ then $G$ is virtually nilpotent.
In all cases, $G$ has derived length at most $3$.
\end{theorem}

\begin{proof}
If $h(\sqrt{G})=1$ then $\sqrt{G}$ is isomorphic to a subgroup of $\mathbb{Q}$ 
and (as in Theorem \ref{h=2}) $G/\sqrt{G}$ has no locally finite normal subgroup.
Since $C_G(\sqrt{G})$ is virtually solvable,
it follows that $C_G(\sqrt{G})=\sqrt{G}$ and so 
$G/\sqrt{G}$ embeds in $Aut(\sqrt{G})$, 
which is isomorphic to a subgroup of $\mathbb{Q}^\times$.
Hence $G/\sqrt{G}\cong\mathbb{Z}^2$,
and so $G$ has derived length 2.

If $h(\sqrt{G})=2$ then $\sqrt{G}$ is abelian and (as in Theorem \ref{h=2} again)
the maximal locally finite normal subgroup of $G/\sqrt{G}$ has order at most 2.
Since $G/\sqrt{G}$ is virtually solvable  and $h(G/\sqrt{G})=1$,
it has an abelian normal subgroup $A$ of rank 1,
which we may assume torsion-free and of finite index in $G/\sqrt{G}$.
Moreover, $G/\sqrt{G}$ embeds in $Aut(\sqrt{G})$, 
which is now isomorphic to a subgroup of $GL(2,\mathbb{Q})$.
No nontrivial element of $A$ can have both eigenvalues roots of unity,
for otherwise $C_G(\sqrt{G})>\sqrt{G}$.
Since the eigenvalues of $A$ have degree $\leq2$ over $\mathbb{Q}$,
it follows that no nontrivial element of $A$ can be infinitely divisible in $A$.
Hence $G/\sqrt{G}$ is virtually $\mathbb{Z}$, 
and so it is either $\mathbb{Z}$ or the infinite dihedral group 
$D_\infty=\mathbb{Z}/2\mathbb{Z}*\mathbb{Z}/2\mathbb{Z}$,
or an extension of one of these by $\mathbb{Z}/2\mathbb{Z}$.

If $G$ has a normal subgroup $H$ such that 
$H/\sqrt{G}\cong\mathbb{Z}/2\mathbb{Z}$ then conjugation in $G$ 
must preserve the filtration $0<H'<\sqrt{G}$ of $\sqrt{G}$.
Therefore elements of $G'$ act nilpotently on $\sqrt{G}$,
and so $G/H$ cannot be $D_\infty$.
Thus if $h(\sqrt{G})=2$ then $G/\sqrt{G}\cong\mathbb{Z}$, 
$D_\infty$ or $\mathbb{Z}\oplus\mathbb{Z}/2\mathbb{Z}$,
and $G$ has derived length 2, 3 or 2, respectively.

If $h(\sqrt{G})=3$ then $h(G/\sqrt{G})=0$, and so $G$ is virtually nilpotent.
Since iterated commutators live in finitely generated subgroups, 
the derived length of $G$ is the maximum of the derived lengths of its finitely generated subgroups.
Finitely generated torsion-free virtually nilpotent groups of Hirsch length 3
are polycyclic, and are fundamental groups of $\mathbb{N}il^3$-manifolds.
These are Seifert fibred over flat 2-orbifolds without reflector curves,
and so these groups have derived length $\leq3$.
Hence $G$ has derived length $\leq3$.
\end{proof}

\end{document}
